\chapter{M\'odulo de Notificaciones}

\section{Descripci\'on General}

El sistema de notificaciones proporciona alertas en tiempo real para
eventos importantes del sistema. Las notificaciones se almacenan en la
base de datos y se entregan al usuario mediante polling peri\'odico.

\section{Esquema}

La tabla \texttt{notifications} almacena:

\begin{itemize}[leftmargin=*]
    \item \texttt{user\_id}: Usuario destinatario.
    \item \texttt{title}: T\'{\i}tulo corto.
    \item \texttt{message}: Mensaje descriptivo.
    \item \texttt{type}: Categor\'{\i}a (appointment, invoice, client,
          patient, system).
    \item \texttt{link}: URL opcional de destino al hacer clic.
    \item \texttt{read}: Estado de lectura (boolean).
\end{itemize}

\section{Server Actions}

\begin{lstlisting}[style=typescript, caption=Notificaciones - Acciones principales]
// Crear notificacion individual
export async function createNotification(data: CreateNotificationInput) {
  await db.insert(schema.notifications).values({
    userId: data.userId,
    title: data.title,
    message: data.message,
    type: data.type,
    link: data.link || null,
  });
  revalidatePath('/');
}

// Obtener notificaciones del usuario (ultimas 50)
export async function getNotifications() {
  const session = await auth();
  if (!session?.user?.id) return [];

  return db
    .select()
    .from(schema.notifications)
    .where(eq(schema.notifications.userId,
              Number(session.user.id)))
    .orderBy(desc(schema.notifications.createdAt))
    .limit(50);
}

// Contar no leidas
export async function getUnreadCount(): Promise<number> {
  const session = await auth();
  if (!session?.user?.id) return 0;

  return db
    .select({ count: sql<number>`count(*)` })
    .from(schema.notifications)
    .where(and(
      eq(schema.notifications.userId, Number(session.user.id)),
      eq(schema.notifications.read, false),
    ))
    .then(r => Number(r[0]?.count ?? 0));
}

// Marcar como leida
export async function markAsRead(id: number) {
  const session = await auth();
  await db
    .update(schema.notifications)
    .set({ read: true })
    .where(and(
      eq(schema.notifications.id, id),
      eq(schema.notifications.userId, Number(session.user.id)),
    ));
  revalidatePath('/notifications');
}
\end{lstlisting}

\section{Notificaciones Masivas}

\subsection{Nueva Cita}

\begin{lstlisting}[style=typescript, caption=notifyNewAppointment]
export async function notifyNewAppointment(data: {
  clientName: string;
  patientName: string;
  date: string;
  time: string;
  vetIds: number[];
}) {
  const title = 'Nueva cita programada';
  const message = `${data.clientName} agendo cita para ` +
                  `${data.patientName} el ${data.date} ` +
                  `a las ${data.time}`;

  // Notificar al veterinario y a todos los administradores
  const admins = await db
    .select({ id: schema.users.id })
    .from(schema.users)
    .where(eq(schema.users.role, 'admin'));

  const allIds = [...new Set([
    ...data.vetIds,
    ...admins.map(a => a.id)
  ])];

  // Insercion bulk
  if (allIds.length > 0) {
    await db.insert(schema.notifications).values(
      allIds.map(userId => ({
        userId, title, message,
        type: 'appointment' as const,
        link: '/appointments',
      }))
    );
  }
}
\end{lstlisting}

\section{Navbar con Campana}

El componente \texttt{Navbar} integra las notificaciones con actualizaci\'on
autom\'atica cada 15 segundos:

\begin{lstlisting}[style=typescript, caption=Navbar - Polling de notificaciones]
// Estado local
const [notifications, setNotifications] = useState<Notification[]>([]);
const [unreadCount, setUnreadCount] = useState(0);

// Fetch inicial y polling cada 15 segundos
const fetchNotifications = useCallback(async () => {
  try {
    const res = await fetch('/api/notifications');
    if (res.ok) {
      const data = await res.json();
      setNotifications(data.notifications ?? []);
      setUnreadCount(data.unread ?? 0);
    }
  } catch {}
}, []);

useEffect(() => {
  fetchNotifications();
  const interval = setInterval(fetchNotifications, 15000);
  return () => clearInterval(interval);
}, [fetchNotifications]);
\end{lstlisting}

\section{Disparadores Autom\'aticos}

Las notificaciones se crean autom\'aticamente en:

\begin{itemize}[leftmargin=*]
    \item \textbf{Crear cita:} Notifica al veterinario asignado y
          administradores.
    \item \textbf{Crear cliente:} Notifica a administradores.
    \item \textbf{Crear factura:} Notifica a administradores y
          recepcionistas.
\end{itemize}

\section{P\'agina de Notificaciones}

La p\'agina \texttt{/notifications} ofrece:

\begin{itemize}[leftmargin=*]
    \item Lista completa de notificaciones.
    \item Icono por tipo (appointment=Calendar, invoice=Receipt,
          client=Users, patient=PawPrint, system=Info).
    \item Resaltado visual de no le\'{\i}das (borde + fondo tenue).
    \item Bot\'on ``Ver'' para notificaciones con enlace.
    \item Botones ``Le\'{\i}do'' individual y ``Marcar todas le\'{\i}das''.
    \item Estado vac\'{\i}o con icono y mensaje.
\end{itemize}

\section{API Endpoint}

\begin{lstlisting}[style=typescript, caption=GET /api/notifications - Endpoint JSON]
export async function GET() {
  const session = await auth();
  if (!session?.user?.id) {
    return NextResponse.json(
      { error: 'No autorizado' }, { status: 401 }
    );
  }

  const [unread, recent] = await Promise.all([
    db.select({ count: sql<number>`count(*)` })
      .from(schema.notifications)
      .where(and(
        eq(schema.notifications.userId, Number(session.user.id)),
        eq(schema.notifications.read, false),
      ))
      .then(r => Number(r[0]?.count ?? 0)),

    db.select()
      .from(schema.notifications)
      .where(eq(schema.notifications.userId,
                Number(session.user.id)))
      .orderBy(desc(schema.notifications.createdAt))
      .limit(5),
  ]);

  return NextResponse.json({ unread, notifications: recent });
}
\end{lstlisting}
